% simulant — Cover letter (Journal of Hazardous Materials)
% 재집필 2026-10-13, Step 3 롤백 방향 C. em-dash 금지 / AI 문체 금지 규칙 적용.
\documentclass[11pt,a4paper]{article}
\usepackage[margin=2.5cm]{geometry}
\usepackage{times}
\usepackage[hidelinks]{hyperref}
\setlength{\parindent}{0pt}
\setlength{\parskip}{8pt}
\pagestyle{empty}

\begin{document}

{\small
Chemical, Biological, Radiological and Nuclear (CBRN) Defense Research Institute\\
Ministry of National Defense, Seoul, Republic of Korea\\
13 October 2026}

\vspace{6pt}

Dear Editors,
% [확인 필요] 담당 에디터 성함 확인 후 "Dear Dr. <Name>," 로 교체.
%            JHM 은 주제별 복수 Editor-in-Chief 체제이므로 투고 시스템에서
%            배정된 에디터를 확인하고 기입할 것.

We submit \textit{Whether Chemical Warfare Agent Simulant Designations Look
Supported Depends on Which Compounds Are in the Comparison} for consideration as
a Research Paper.

The manuscript tests a convention this journal recently reviewed. Webb, Coulon
and Temple (\textit{J.\ Hazard.\ Mater.} \textbf{492}, 138021, 2025) state the
criterion by which low-toxicity chemical warfare agent simulants are chosen:
align hydrophobicity, volatility and solubility with the agent, and the compounds
so chosen present minimal toxicity risk. Both halves are testable against
existing measurement. We tested them. The toxicity half gave a clear answer. The
fidelity half gave an answer that depends on what the chosen compounds are
compared against, and that is the paper.

\textbf{On a panel that contains the compounds the field itself uses as
simulants, none of eleven agents has a valid match to an incumbent simulant.}
For eleven authentic agents we computed RDKit structural and GFN2-xTB electronic
descriptors and asked, by exact permutation over every possible designation,
whether dimethyl methylphosphonate, dichloromethane and dipropylene glycol
monomethyl ether sit closer to the agents than compounds drawn at random. Our
primary panel of thirty-five compounds includes five further compounds the
sensor, sorbent and decontamination literature uses as simulants (DFP, DCP,
DCNP, CEES, methyl salicylate). On it, not one agent is matched, and not one of
eighteen tests in three descriptor spaces under six distance statistics reaches
nominal significance. The mechanism is not a loss of fidelity. The incumbent compounds sit
almost exactly where they sat before: sarin moves from 1.14 to 1.19 typical
nearest-neighbour separations from DMMP. What the larger panel supplies is
diisopropyl fluorophosphate at 1.07 from sarin, 0.99 from soman and 1.17 from
cyclosarin, and 2-chloroethyl ethyl sulfide at 0.57 from sulfur mustard against
dichloromethane's 1.67. Resampling the 41 descriptors 1000 times leaves this
intact: the incumbent three match a median of zero agents (95\% interval zero to
three) against five (four to six) for the eight compounds the field uses, and
the difference is positive in every replicate. A designation defended by a
comparison that excludes the alternatives is not defended.

\textbf{The twenty-eight-compound panel that omits those five is reported in
full, because the contrast is the result.} On it, three agents are matched, DMMP
supplies all three, and the smallest nominal $p$ is 0.025. How that is corrected
decides whether it means anything, so we report five schemes rather than one:
under Benjamini-Hochberg across the eighteen tests, across six within each
space, at the effective number of independent tests estimated from the
correlation of the permutation nulls, or under the Westfall-Young min-$P$
procedure, which handles the dependence among the eighteen exactly because they
share one enumeration of designations, nothing reaches 0.05 (the Westfall-Young
value, 0.075, is the closest any correction brings it). Had the statistic
been specified in advance, it would survive at $q=0.042$. We did not specify one.
That contrast is the paper's argument made with its own data. The observed value
also sits at the 25.8th percentile of 2000 arbitrary three-descriptor subspaces
of the same descriptor block.

\textbf{The three incumbent simulants were designated for measurements that are
not the one they now support.} We looked for the instrument that designates
them. DMMP was identified as a filtration-breakthrough test simulant in a NIOSH
respirator programme, in a document that also declines to recommend any
simulant. DPGME appears as ``a preferred G agent simulant'' in a US Army patent,
in a training composition carrying a fluorescent dye and justified on
ion-mobility detectability. For dichloromethane we found no designating document
at all; it is absent from the NIOSH candidate list and from the Man-in-Simulant
Test programme. The compounds with the clearest designations are not these
three: methyl salicylate is the Army's Man-in-Simulant Test simulant for sulfur
mustard, and the NIOSH candidate list is DFP, DIMP, DEEP, TEP, dibutyl sulfide,
CEES, malathion and methyl salicylate. A simulant is validated against a
measurement, not against an agent, and that validation does not travel with the
compound into sensor development.

\textbf{Consistent with that, we cannot establish that the incumbent simulants
resemble the agents, and we cannot establish that they do not.} Three things are
open to the analyst: the compound panel, the descriptor space, and the distance
statistic. We fixed all three in every combination we could compute, which is
eighteen tests on each of three panel and designation configurations, and we
report the whole table rather than a cell from it.

\textbf{One widely used descriptor turned out not to be a coordinate.} The Parr
electrophilicity index $\omega = \mu^{2}/2\eta$ is an exact function of the two
frontier orbital energies normally reported beside it; across our panel the
deviation from the definition is at most $3.6\times10^{-15}$~eV. A correlation
filter does not catch this, because $r(E_\mathrm{LUMO}, \omega) = -0.899$. We
removed it, which cost the electronic space a dimension it never had, and we
state the general form of the point in the Discussion because any descriptor set
reporting $E_\mathrm{HOMO}$, $E_\mathrm{LUMO}$, the gap, $\mu$, $\eta$ and
$\omega$ as six quantities has two.

\textbf{The measured array agrees that the question is unresolved, and corrects
us on one point.} Twenty-nine-dye colorimetric response for twelve of the
compounds, tested by Mantel permutation over compounds rather than over the
pairs they generate, gives two nominally significant cells out of eight that
favour different descriptor spaces, neither of which survives correction across
the eight (smallest $q$ = 0.098). Two artifacts of this kind of array analysis are
reported alongside, because our own first version contained both: dye-axis
centring fixes the between-compound mean cosine at $-1/(n-1)$, and our observed
means clear that floor by only 0.013 to 0.066, so a negative between-compound
similarity is arithmetic rather than evidence; and the measured agent nearest to
DMMP in the ultraviolet channel is tabun, not sarin, although tabun fails our own
descriptor matching criterion.

\textbf{On safety the two incumbents differ in evidence, not only in hazard.}
Dichloromethane has a regulator-derived acute inhalation value and it is
favourable on that axis: at 41\,685~mg~m$^{-3}$ it sits near the safe end of the
nineteen panel compounds that carry one, though methanol and n-hexane score
higher still. We are explicit in the manuscript that this is the acute axis only
and says nothing about the chronic and carcinogenic classification that is
actually displacing dichloromethane from solvent inventories. DMMP has no value
at all, because EPA declined to derive an inhalation reference concentration for
want of data. Reporting the two as one category of low-toxicity simulant
obscures that difference, which matters because DMMP is also the compound against
which decontamination efficiency is quantified.

\textbf{A model trained on the regulatory inventory cannot fill the gap, and the
failure is instructive.} A random forest on 230 AEGL chemicals reaches
$R^{2}=0.458$ with confidence intervals that do not overlap its baseline. Applied
to the twenty military agents held out of training, it predicts every one of them
safer than measured, by factors of 2.5 to 2600, and certifies fourteen of the
twenty as inside its applicability domain. We checked the obvious explanation and
it is wrong: thirteen training compounds carry an acetylcholinesterase-inhibiting
pharmacophore, identified by substructure and listed in the paper, and one of
them is more acutely toxic than sarin. The model fails because it regresses both
tails toward the mean, a bias defined in the response space that an input-space
applicability domain cannot see.

\textbf{The two axes are independent.} Across 207 compounds carrying both
measured axes, Spearman $r = +0.004$ (95\% CI $-0.132$ to $+0.143$). Eight
compounds combine top-quartile fidelity with measured acute inhalation safety
within an order of magnitude of dichloromethane, all eight clear a structural
filter for agent-class motifs, and seven of the eight are more faithful than
dichloromethane itself. We propose none of them as a replacement, and the reason
is the paper's own argument: a recommendation resting on one panel, one space and
two axes would repeat the error we are reporting.

The work sits within the scope of this journal on both sides. The simulant
selection literature it tests is published here, and the hazard-screening
methods it uses, regulator-derived acute toxicity values bounded by an
applicability domain, are the methods this journal applies to other compound
classes.

\textbf{On data integrity, and on this manuscript's own history.} This analysis
was audited four times before submission and each audit changed the numbers, so
we state what happened rather than present the final version as though it had
come out that way.

The first audit found that the numeric parser had split thousands separators in
the source table, so that ``12,000~ppm'' was read as 12. The second, of chemical
identity rather than of numbers, found two records whose stored structure
belonged to a different compound; both are excluded. The third found four further
faults in how the regulatory table was read, all of one kind: treating a printed
table as though its cells always held plain numbers. Footnote lines beginning
\textit{AEGL 3 - 10 min = ...} were parsed as though they were table rows, which
turned duration labels into values and gave propane and butane ten-minute
AEGL-3 values of 10~ppm against the true 33,000 and 77,000~ppm. Nine compounds
whose ten-minute cell reads \textit{see below}, the agency's marker for a value
above half the lower explosive limit, had no value collected at all, among them
methanol, acetone, benzene and n-hexane. Rows whose longer durations read
\textit{NR} were discarded whole, which silently dropped three arsenical
vesicants and BZ from the held-out agent set. And one row broken across two lines
by the document's text layer was lost.

The fourth audit was of our own statistics rather than of the data, and it is the
reason the paper now says what it says. It found the functional dependence in the
electronic descriptor set, the pairwise-distance pseudoreplication in the array
analysis, the algebraic floor created by centring, the absence of any multiplicity
correction across eighteen tests we had in fact computed, and the fact that we had
never run the panel that includes the simulants the field uses. An earlier version
of this manuscript reported $p=0.028$ as a result and built a thesis on it. None
of those five problems is visible in that version, and together they remove the
thesis. We have rewritten the paper around what the data support, and we report
the earlier reading alongside it rather than deleting it, because how a nominal
result of this kind arises is part of what a reader of this literature needs.

The extraction is now cell-based and refuses any row that does not resolve into
five duration cells, and it is validated over the whole table by a second reader
that shares neither the text-extraction engine nor the algorithm, reading values
by character column position rather than by token order. The two readers agree on
all 259 entries, on all 257 of those that resolved to a structure and on all
nineteen panel values.

\textbf{Dual-use.} The optimisation direction is fixed throughout and stated in
the manuscript: detection-relevant fidelity is maximised while phosphorus
electrophilicity, acetylcholinesterase reactivity, alkylating capacity and
predicted mammalian toxicity are minimised. The structural filter excludes
candidates and never orders them. We did not extend the toxicity model to the
organophosphorus class, because a predictor of mammalian lethality evaluated over
structures not yet made, inside a loop that proposes them, is the configuration
this work is built to avoid. We cite the existing literature on organophosphorus
lethality modelling, including our own, rather than implying that none exists.

The acute toxicity axis, the random-forest interpolation of it, the decoupling
test between fidelity and toxicity, and the phosphorus-local reactivity analysis
are in the Supplementary Material. They stand on their own data and are
summarised where the main text needs them; we placed them there because the
argument of the paper turns on the composition of the comparison panel and does
not depend on them.

The manuscript is original, is not under consideration elsewhere, and has been
approved by both authors. The authors declare no competing financial interest,
and the work received no specific funding. Analysis code, descriptor tables,
parsed toxicity values with their sources and all figure scripts are deposited in
the public repository cited in the manuscript and also accompany the submission.
The exact permutation tests are exhaustive and the remaining analyses are
deterministic given the stated seeds.

Thank you for considering this work.

\vspace{10pt}

Yours sincerely,

\vspace{14pt}

Yoon Jeong Jang, corresponding author\\
on behalf of Ku Kang and Yoon Jeong Jang\\
\texttt{yjjangchem@korea.kr}

% 추천 심사위원: Ku 결정에 따라 제출 시스템에서 공란으로 둔다 (2026-10-15).

\end{document}
